\chapter{A First Program}\label{ch-first-program}

\section{Hello}

Make a file called \verb|hello.zeno| with one line in it:

\begin{zeno}[file=hello.zeno, cmd=zeno run hello.zeno]
print("Hello, Zeno!")
\end{zeno}

That is the whole program. There is no \verb|main| function to write,
nothing to import and no semicolon at the end. Zeno starts at the top
of the file and does each line in turn, and when it runs out of lines
the program is over.

\verb|print| shows what it is given and then ends the line. Give it
several things and it puts a space between each; give it nothing and
it prints an empty line:

\begin{zeno}
print("one", "two", 3)
print()
print("after an empty line")
\end{zeno}

Text is written between double quotes, \verb|"like this"|. Numbers
are written as they are. Each has a chapter of its own later.

\section{Comments}

A \verb|#| starts a comment. Everything from it to the end of the line
is for the reader, and Zeno skips it:

\begin{zeno}
# Greet the world. This whole line is a comment.
print("Hello")   # and so is the end of this one
\end{zeno}

\section{One line, one statement}

Each line is one \emph{statement}: one thing for the program to do.
The end of the line ends the statement, so Zeno needs no \verb|;|.

A long statement can go on to the next line when its first line is
clearly not finished. That is so inside brackets, and after anything
that cannot end a statement, such as \verb|+|, \verb|,| or \verb|=|:

\begin{zeno}
let total = 1 +
    2 +
    3
print(total,
      "is the total")
\end{zeno}

The other way round does not work. A line that could stand on its own
\emph{does} stand on its own, so a line that starts with \verb|+| is a
mistake:

\begin{zeno}[fails=yes]
let total = 1
    + 2
\end{zeno}

\begin{why}
The line break already means "next statement" to the person reading,
so a \verb|;| on every line only says it again. The rule for going on
to the next line, "an unfinished line continues", can never be misread,
because a line that looks finished is finished.
\end{why}

\section{Names for things}

\verb|let| gives a value a name, and from then on the name stands for
the value:

\begin{zeno}
let name = "Ada"
let year = 1843
print(name, "wrote the first program in", year)
print(f"Hello, {name}!")
\end{zeno}

The last line is a \emph{formatted string}. The \verb|f| in front of
the quotes says that what stands inside \verb|{ }| is to be worked out
and put into the text. \hyperref[ch-values]{Values and Variables} and
\hyperref[ch-text]{Text} say much more about both.

\section{When something is wrong}

Zeno reads the whole program before it runs any of it, and if it finds
a mistake it says so and runs nothing. Here is a program with two:

\begin{zeno}[fails=yes]
let name = "Ada"
name = "Grace"
print(nmae)
\end{zeno}

Each message says what is wrong, and where: the file, the line and the
column, the line itself, and a \verb|^| under the place. Where Zeno can
tell what you meant, it says what to write instead. The first mistake
here is that a name given with \verb|let| cannot be given another
value, and the message names the word that would allow it. The second
is a misspelt name.

Habits from other languages get messages of their own:

\begin{zeno}[fails=yes]
def greet():
    print("hi")
\end{zeno}

\begin{zeno}[fails=yes]
let ready = true
if ready && true {
    print("go")
}
\end{zeno}

\begin{zeno}[fails=yes]
print("hello");
\end{zeno}

\begin{why}
An error message is the language teaching the person who reads it,
at the moment the lesson is most useful. So every message in Zeno
aims to say what to write instead, not only what is wrong.
\end{why}

Notice that the last three programs each showed one message, while the
one before showed two. A mistake in the \emph{shape} of a program,
such as a missing brace, leaves Zeno unsure where it is, and what it
would report after that could be nonsense caused by the first. So it
stops at the first. A mistake in the \emph{meaning}, such as a
misspelt name, leaves it sure where it is, so it reports every one it
finds.

\section{Checking without running}

\verb|zeno run| checks a program and then runs it. \verb|zeno check|
only checks it:

\begin{zeno}[file=greet.zeno, cmd=zeno check greet.zeno]
let name = "Ada"
print(f"Hello, {name}!")
\end{zeno}

It prints nothing when there is nothing wrong. That makes it useful
for a program that would take a long time to run, or one that drives
hardware you do not want to move yet.

\begin{summary}
A program is a file of statements, done from the top, one to a line.
\verb|print| shows values, \verb|#| starts a comment, and \verb|let|
names a value. Zeno checks the whole program before it runs a line of
it, and a message says what is wrong, where, and usually what to write
instead.
\end{summary}

\begin{tryit}
\begin{enumerate}
\item Print your own name and the year you were born on one line,
  with one \verb|print|.
\item Leave the closing quote off a string and read the message. Then
  leave off a closing bracket.
\item Write \verb|elif| where Zeno expects \verb|else if|, and
  \verb|x := 1| for a variable. What does Zeno say to each?
\item Put two mistakes of meaning in one program and run
  \verb|zeno check|. How many messages are there? Now add a mistake of
  shape, such as \verb|print("a"|. How many now, and why?
\end{enumerate}
\end{tryit}
